% !TeX root = ../colloquium-2.tex

% Формулировки первого и второго правил Лопиталя.

\phantomsection\label{sec_Definition14}
\begin{definition}[Формулировки первого и второго правил Лопиталя.]
    \subsection*{Формулировка первого правила Лопиталя, 
                 неопределённость вида $\frac{0}{0}$.}
    \begin{note}[purple]
        Пусть:
        \begin{enumerate}
            \item Функции $f$ и $g$ определены и дифференцируемы на интервале $(a, \; b)$
            \item $ \ds \lim_{x \rightarrow b-} f(x) = \lim_{x \rightarrow b-} g(x) = 0 $
            \item $g'(\xi) \neq 0 \; \forall \xi \in (a, \; b)$
            \item Существует конечный или бесконечный предел 
            $\ds \lim_{x \rightarrow b-} \frac{f'(x)}{g'(x)}$
        \end{enumerate}
        Тогда справедливо равенство 
        $\ds \lim_{x \rightarrow b-} \frac{f(x) }{g(x) } = 
                       \lim_{x \rightarrow b-} \frac{f'(x)}{g'(x)} $
    \end{note}

    \subsection*{Формулировка второго правила Лопиталя,
                 неопределённость вида $\frac{\infty}{\infty}$.}
    \begin{note}[purple]
        \begin{enumerate}
            \item Функции $f$ и $g$ определены и дифференцируемы на интервале $(a, \; b)$
            \item $\lim_{x \rightarrow b-} g(x) = \infty$
            \item $g'(\xi) \neq 0 \; \forall \xi \in (a, \; b)$
            \item Существует конечный или бесконечный предел 
            $\ds \lim_{x \rightarrow b-} \frac{f'(x)}{g'(x)}$
        \end{enumerate}
        Тогда справедливо равенство $\ds \lim_{x \rightarrow b-} \frac{f(x)} {g(x)} = 
                                         \lim_{x \rightarrow b-} \frac{f'(x)}{g'(x)}$
    \end{note}
\end{definition}